\documentclass[12pt]{article}
\usepackage{amsmath, amssymb}
\usepackage[margin=1in]{geometry}
\setlength{\parskip}{6pt}
\title{QUBO Formulation: Molecular Conformation Selection Problem}
\date{}
\begin{document}
\maketitle

\subsection*{Molecular Conformation Selection Problem}

\paragraph{Problem Statement.}
Given a molecule decomposed into $N_t$ rotatable torsions, each torsion $t \in \{0, \dots, N_t-1\}$ admits a discrete set of possible rotamer states $r \in \{0, \dots, N_r(t)-1\}$. Each rotamer configuration carries a local energy $E_{t,r}$, and every pair of rotamers at torsions $t_1$ and $t_2$ contributes an interaction energy $I_{t_1, t_2, r_1, r_2}$. The task is to select exactly one rotamer per torsion such that the total energy, defined as the sum of local energies of the selected rotamers plus the sum of interaction energies between all selected pairs, is minimized.

\paragraph{Decision Variables.}
For each torsion $t \in \{0, \dots, N_t-1\}$ and each rotamer $r \in \{0, \dots, N_r(t)-1\}$, introduce a binary decision variable $x_{t,r} \in \{0,1\}$. The variable $x_{t,r}$ equals $1$ if rotamer $r$ is selected for torsion $t$, and $0$ otherwise.

\paragraph{QUBO Derivation.}
Step 1 --- Write the original objective in general notation. The minimization objective is
\begin{align}
H_{\text{obj}}(x) = \sum_{t=0}^{N_t-1} \sum_{r=0}^{N_r(t)-1} E_{t,r} x_{t,r} + \sum_{t_1=0}^{N_t-1} \sum_{r_1=0}^{N_r(t_1)-1} \sum_{t_2=0}^{N_t-1} \sum_{r_2=0}^{N_r(t_2)-1} I_{t_1, t_2, r_1, r_2} x_{t_1, r_1} x_{t_2, r_2}.
\end{align}

Step 2 --- Write each constraint in its original algebraic form. For each torsion $t$, exactly one rotamer must be selected:
\begin{align}
\sum_{r=0}^{N_r(t)-1} x_{t,r} = 1, \quad \forall t \in \{0, \dots, N_t-1\}.
\end{align}

Step 3 --- Convert each constraint to a penalty term. Let $\lambda > 0$ be a penalty weight. For a fixed torsion $t$, the squared penalty is
\begin{align}
P_t(x) = \lambda \left( \sum_{r=0}^{N_r(t)-1} x_{t,r} - 1 \right)^2.
\end{align}
Expanding the square yields
\begin{align}
P_t(x) = \lambda \left( \sum_{r=0}^{N_r(t)-1} x_{t,r}^2 + \sum_{\substack{r, r'=0 \\ r \neq r'}}^{N_r(t)-1} x_{t,r} x_{t,r'} - 2 \sum_{r=0}^{N_r(t)-1} x_{t,r} + 1 \right).
\end{align}
Applying the idempotent property $x_{t,r}^2 = x_{t,r}$ for binary variables simplifies the expression to
\begin{align}
P_t(x) = \lambda \left( \sum_{r=0}^{N_r(t)-1} x_{t,r} + \sum_{\substack{r, r'=0 \\ r \neq r'}}^{N_r(t)-1} x_{t,r} x_{t,r'} - 2 \sum_{r=0}^{N_r(t)-1} x_{t,r} + 1 \right) = \lambda \left( \sum_{\substack{r, r'=0 \\ r \neq r'}}^{N_r(t)-1} x_{t,r} x_{t,r'} - \sum_{r=0}^{N_r(t)-1} x_{t,r} + 1 \right).
\end{align}
The resulting general penalty contribution to the Hamiltonian is $\sum_{t=0}^{N_t-1} P_t(x)$.

Step 4 --- Combine objective and all penalty terms into the single QUBO Hamiltonian $H(x)$:
\begin{align}
H(x) = \sum_{t=0}^{N_t-1} \sum_{r=0}^{N_r(t)-1} E_{t,r} x_{t,r} + \sum_{t_1, t_2, r_1, r_2} I_{t_1, t_2, r_1, r_2} x_{t_1, r_1} x_{t_2, r_2} + \sum_{t=0}^{N_t-1} \lambda \left( \sum_{\substack{r, r'=0 \\ r \neq r'}}^{N_r(t)-1} x_{t,r} x_{t,r'} - \sum_{r=0}^{N_r(t)-1} x_{t,r} + 1 \right).
\end{align}

Step 5 --- Read off the Q matrix entries and offset $c$ from $H$ in general closed form. Grouping linear and quadratic terms yields the standard QUBO form $H(x) = \sum_{(t,r)} Q_{(t,r),(t,r)} x_{t,r} + \sum_{(t,r) < (t',r')} Q_{(t,r),(t',r')} x_{t,r} x_{t',r'} + c$, where the entries are:
\begin{align}
Q_{(t,r),(t,r)} &= E_{t,r} - \lambda, \\
Q_{(t,r),(t,r')} &= \lambda, \quad \text{for } r \neq r', \\
Q_{(t_1,r_1),(t_2,r_2)} &= I_{t_1, t_2, r_1, r_2}, \quad \text{for } t_1 \neq t_2, \\
c &= N_t \lambda.
\end{align}
By symmetry, $Q_{(t,r),(t',r')} = Q_{(t',r'),(t,r)}$.

\paragraph{Penalty Weight Justification.}
Let $W_{\max} = \max_{t,r} |E_{t,r}| + \sum_{t', r'} |I_{t,r,t',r'}|$. Choosing $\lambda > W_{\max}$ guarantees that any violation of the constraint $\sum_r x_{t,r} = 1$ incurs a penalty of at least $\lambda$, which strictly exceeds the maximum possible reduction in the objective function achievable by selecting zero or multiple rotamers for torsion $t$. Consequently, the energy of any infeasible assignment is strictly greater than that of the optimal feasible assignment.

\paragraph{Correctness Argument.}
Minimizing $H(x)$ over $\{0,1\}^n$ recovers the optimal molecular conformation because the penalty terms vanish if and only if all constraints are satisfied. For feasible assignments, $H(x)$ exactly equals the original molecular energy. For infeasible assignments, the penalty contribution is at least $\lambda$, which dominates any potential decrease in the objective function from constraint violations. Therefore, the global minimum of $H(x)$ must occur at a feasible point, and among all feasible points, it corresponds to the assignment that minimizes the original energy.

\paragraph{Illustrative Example.}
Consider an instance with $N_t = 2$ torsions and rotamer counts $[2, 2]$. The binary variables are indexed as $x_0 = x_{0,0}$, $x_1 = x_{0,1}$, $x_2 = x_{1,0}$, and $x_3 = x_{1,1}$. Instantiating the general formulas yields the concrete objective:
\begin{align}
H_{\text{obj}}(x) = E_{0,0}x_0 + E_{0,1}x_1 + E_{1,0}x_2 + E_{1,1}x_3 + I_{0,0,0,1}x_0x_1 + I_{1,1,0,1}x_2x_3 + (I_{0,1,0,0}+I_{1,0,0,0})x_0x_2 + (I_{0,1,0,1}+I_{1,0,0,1})x_0x_3 + (I_{0,1,1,0}+I_{1,0,1,0})x_1x_2 + (I_{0,1,1,1}+I_{1,0,1,1})x_1x_3.
\end{align}
The concrete penalty terms are:
\begin{align}
P_0(x) = \lambda(x_0 + x_1 - 1)^2, \quad P_1(x) = \lambda(x_2 + x_3 - 1)^2.
\end{align}
Expanding and combining with the objective gives the resulting Hamiltonian $H(x) = \sum_{i=0}^{3} Q_{i,i} x_i + \sum_{0 \le i < j \le 3} Q_{i,j} x_i x_j + c$, with the following Q matrix entries:
\begin{align}
Q_{0,0} &= E_{0,0} - \lambda, & Q_{1,1} &= E_{0,1} - \lambda, & Q_{2,2} &= E_{1,0} - \lambda, & Q_{3,3} &= E_{1,1} - \lambda, \\
Q_{0,1} &= \lambda + I_{0,0,0,1}, & Q_{2,3} &= \lambda + I_{1,1,0,1}, & Q_{0,2} &= I_{0,1,0,0} + I_{1,0,0,0}, \\
Q_{0,3} &= I_{0,1,0,1} + I_{1,0,0,1}, & Q_{1,2} &= I_{0,1,1,0} + I_{1,0,1,0}, & Q_{1,3} &= I_{0,1,1,1} + I_{1,0,1,1}, \\
c &= 2\lambda.
\end{align}
For this small instance, the optimal assignment can be found by enumerating the $2^4 = 16$ binary vectors, evaluating $H(x)$, and selecting the $x^*$ that minimizes the energy while satisfying $\sum_{r} x_{t,r} = 1$ for $t \in \{0,1\}$. The minimum value $H(x^*)$ gives the optimal conformation energy.

\end{document}